% AI-translated component: printed 162--171.
\begingroup
\clearpage
\noindent162\hfill\textsc{Posthumous Papers.}\hfill\mbox{}\par\medskip
\begin{center}
\begingroup\fontsize{7.6}{9}\selectfont\setlength{\tabcolsep}{2pt}\renewcommand{\arraystretch}{1.07}
\begin{tabular}{|r|r|r|r|r||r|r|r|r|r|}
\hline
\rule{0pt}{18pt}\shortstack{$x$\\(metres)}&$\varphi$&\shortstack{$\log\alpha$\\$1{,}....-10$}&\shortstack{$\log\beta$\\$7{,}....-30$}&\shortstack{$\log\gamma$\\$2{,}....-20$}&\shortstack{$x$\\(metres)}&$\varphi$&\shortstack{$\log\alpha$\\$1{,}....-10$}&\shortstack{$\log\beta$\\$7{,}....-30$}&\shortstack{$\log\gamma$\\$2{,}....-20$}\\\hline
$-200\,000$&$53^\circ19'40''{,}521$&$1{,}4033800$&$7{,}31422$&$2{,}40145$&$-500\,000$&$56^\circ1'24''{,}643$&$1{,}4031246$&$7{,}31397$&$2{,}38676$\\
$210\,000$&$25\ 4{,}061$&$713$&$22$&$103$&$510\,000$&$6\ 48{,}041$&$163$&$96$&$620$\\
$220\,000$&$30\ 27{,}597$&$627$&$21$&$060$&$520\,000$&$12\ 11{,}434$&$1{,}4031079$&$95$&$563$\\
$230\,000$&$35\ 51{,}127$&$541$&$20$&$2{,}40017$&$530\,000$&$17\ 34{,}822$&$1{,}4030996$&$94$&$506$\\
$240\,000$&$41\ 14{,}653$&$455$&$19$&$2{,}39974$&$540\,000$&$22\ 58{,}205$&$912$&$93$&$448$\\
\hline
$-250\,000$&$46\ 38{,}173$&$369$&$18$&$930$&$-550\,000$&$28\ 21{,}584$&$829$&$92$&$389$\\
$260\,000$&$52\ 1{,}689$&$282$&$17$&$885$&$560\,000$&$33\ 44{,}958$&$746$&$92$&$330$\\
$270\,000$&$53\ 57\ 25{,}200$&$197$&$16$&$840$&$570\,000$&$39\ 8{,}329$&$663$&$91$&$271$\\
$280\,000$&$54\ 2\ 48{,}707$&$111$&$15$&$795$&$580\,000$&$44\ 31{,}693$&$580$&$90$&$211$\\
$290\,000$&$8\ 12{,}208$&$1{,}4033025$&$15$&$749$&$590\,000$&$49\ 55{,}053$&$498$&$\raisebox{.6ex}{\tiny\textquoteleft}89$&$150$\\
\hline
$-300\,000$&$13\ 35{,}705$&$1{,}4032939$&$14$&$703$&$-600\,000$&$56\ 55\ 18{,}409$&$415$&$88$&$089$\\
$310\,000$&$18\ 59{,}197$&$854$&$13$&$656$&$610\,000$&$57\ 0\ 41{,}760$&$333$&$87$&$2{,}38028$\\
$320\,000$&$24\ 22{,}684$&$768$&$12$&$609$&$620\,000$&$6\ 5{,}106$&$250$&$87$&$2{,}37966$\\
$330\,000$&$29\ 46{,}167$&$683$&$11$&$561$&$630\,000$&$11\ 28{,}448$&$168$&$86$&$904$\\
$340\,000$&$35\ 9{,}644$&$598$&$10$&$513$&$640\,000$&$16\ 51{,}785$&$086$&$85$&$841$\\
\hline
$-350\,000$&$40\ 33{,}117$&$512$&$09$&$464$&$-650\,000$&$22\ 15{,}118$&$1{,}4030004$&$84$&$777$\\
$360\,000$&$45\ 56{,}585$&$427$&$09$&$415$&$660\,000$&$27\ 38{,}446$&$1{,}4029922$&$83$&$713$\\
$370\,000$&$51\ 20{,}049$&$342$&$08$&$366$&$670\,000$&$33\ 1{,}770$&$840$&$83$&$648$\\
$380\,000$&$54\ 56\ 43{,}508$&$257$&$07$&$316$&$680\,000$&$38\ 25{,}089$&$759$&$82$&$583$\\
$390\,000$&$55\ 2\ 6{,}961$&$173$&$06$&$265$&$690\,000$&$43\ 48{,}403$&$677$&$81$&$518$\\
\hline
$-400\,000$&$7\ 30{,}411$&$088$&$05$&$214$&$-700\,000$&$49\ 11{,}713$&$596$&$80$&$452$\\
$410\,000$&$12\ 53{,}855$&$1{,}4032003$&$04$&$162$&$710\,000$&$54\ 35{,}018$&$514$&$79$&$385$\\
$420\,000$&$18\ 17{,}295$&$1{,}4031919$&$03$&$110$&$720\,000$&$57\ 59\ 58{,}319$&$433$&$78$&$318$\\
$430\,000$&$23\ 40{,}730$&$834$&$03$&$058$&$730\,000$&$58\ 5\ 21{,}615$&$352$&$78$&$250$\\
$440\,000$&$29\ 4{,}160$&$750$&$02$&$2{,}39005$&$740\,000$&$10\ 44{,}907$&$271$&$77$&$182$\\
\hline
$-450\,000$&$34\ 27{,}586$&$666$&$01$&$2{,}38951$&$-750\,000$&$16\ 8{,}194$&$190$&$76$&$113$\\
$460\,000$&$39\ 51{,}007$&$582$&$7{,}31400$&$897$&$760\,000$&$21\ 31{,}476$&$110$&$75$&$2{,}37044$\\
$470\,000$&$45\ 14{,}423$&$498$&$7{,}31399$&$843$&$770\,000$&$26\ 54{,}755$&$1{,}4029029$&$74$&$2{,}36974$\\
$480\,000$&$50\ 37{,}834$&$414$&$98$&$788$&$780\,000$&$32\ 18{,}028$&$1{,}4028949$&$73$&$904$\\
$490\,000$&$55\ 56\ 1{,}241$&$330$&$98$&$732$&$790\,000$&$37\ 41{,}297$&$869$&$73$&$833$\\
\hline
$-500\,000$&$56\ 1\ 24{,}643$&$1{,}4031246$&$7{,}31397$&$2{,}38676$&$-800\,000$&$58\ 43\ 4{,}562$&$1{,}4028788$&$7{,}31372$&$2{,}36762$\\
\end{tabular}\par\endgroup
\end{center}
[This table is based on the flattening value $\frac1{302{,}68}$. The values of $x$ are to be corrected by $-26{,}264\,\mathrm m$ if $\varphi=51^\circ31'47''{,}85$ is taken for Göttingen instead of $51^\circ31'48''{,}70$.]
\bigskip
\begin{center}
[The Difference Between the Projection of a Geodesic\\and the Straight Line Joining Its Endpoints in the Conformal\\Mapping of a Curved Surface in the Plane.]\par\medskip
[16.]
\end{center}
A very fruitful principle for the general theory of curved surfaces is the following.

The individual points of the curved surface are [represented] by two variable quantities $p$, $q$ in any desired manner. All points
\clearpage
\noindent\mbox{}\hfill\textsc{Conformal Mapping of the Spheroid in the Plane.}\hfill163\par\medskip
\noindent for which $q$ is constant thus form a line on the surface. Let its element be $\mathrm du=t\,\mathrm dp$, where $t$ will be a function of $p$ and $q$.

Now draw a shortest line to the arbitrary point of the curved surface determined by $p$, $q$, either from a fixed point or perpendicular to a given line. Its length $r$ will thus itself be a function of $p$ and $q$. If $\varphi$ is the angle that the element of $r$ makes with the element $\mathrm du$, then
\[
\frac{\partial r}{\partial p}=t\cos\varphi.
\]
In the plane representation that preserves the similarity of infinitesimal elements, let a shortest line on the curved surface be represented by the curve whose points are determined by $R$, the distance from the starting point, [and] $\theta$, the direction. Let $n^*$ be the [reciprocal] scale factor, i.e. corresponding elements in the plane and on the surface are in the ratio $1$ to $n^*$.

Further, let
\[
\begin{aligned}
n^*&=N+aR+bRR+cR^3+\ldots\\
\rlap{\hspace{-1.4em}\raisebox{.7ex}{.}}\log n^*&=\log N+\alpha R+\beta RR+\gamma R^3+\ldots,
\end{aligned}
\]
where $a$, $b$, $c$, $\ldots$ and $\alpha$, $\beta$, $\gamma$, $\ldots$ denote functions of $\theta$.

We then have [$\mathrm dn^*=n^*\mathrm d\log n^*$ or]
\[
(a+2bR+3cRR\ldots)=(N+aR+bRR\ldots)(\alpha+2\beta R+3\gamma RR\ldots)
\]
[and hence]
\[
\begin{aligned}
a&=N\alpha\\
2b&=2N\beta+a\alpha\\
3c&=3N\gamma+2a\beta+b\alpha
\end{aligned}
\]
\begin{center}etc.\end{center}
[or]
\[
\begin{aligned}
a&=N\alpha\\
b&=N\beta+\tfrac12N\alpha\alpha\\
c&=N\gamma+N\alpha\beta+\tfrac16N\alpha^3
\end{aligned}
\]
\begin{center}etc.\end{center}
\vfill\hfill$21^*$\hspace*{2cm}\par
\clearpage
\noindent164\hfill\textsc{Posthumous Papers.}\hfill\mbox{}\par\medskip
\begingroup\small
From the above we now have
\[
\begin{gathered}
\frac{\partial r}{\partial R}=n^*\cos\varphi\\
\frac{\partial r}{R\partial\theta}=-n^*\sin\varphi,\\
\left(\frac{\partial r}{\partial R}\right)^2+\left(\frac{\partial r}{R\partial\theta}\right)^2={n^*}^2;
\end{gathered}
\]
[furthermore, we have
\[
m\frac{\partial(\theta+\psi)}{\partial R}=n^*\sin\varphi,\qquad m\frac{\partial(\theta+\psi)}{R\partial\theta}=n^*\cos\varphi,
\]
where $m\,\mathrm d(\theta+\psi)$ denotes the element of the curve of constant $r$. From this it follows that
\[
\frac{\partial(\theta+\psi)}{\partial R}\cos\varphi-\frac{\partial(\theta+\psi)}{R\partial\theta}\sin\varphi=0
\]
or]
\[
\frac{\partial(\theta+\psi)}{\partial R}.\frac{\partial r}{\partial R}+\frac{\partial(\theta+\psi)}{R\partial\theta}.\frac{\partial r}{R\partial\theta}=0.
\]
[$\psi$ and $\varphi$ denote the angles that the image of the geodesic in the plane makes with $R$ at the starting point and endpoint, respectively.]

From this one finds, on setting
\[
\frac{\partial a}{\partial\theta}=a',\qquad\frac{\partial b}{\partial\theta}=b',\ \text{etc.;}\qquad
\frac{\partial a'}{\partial\theta}=a'',\qquad\frac{\partial b'}{\partial\theta}=b'',\ \text{etc.}
\]
and noting that $a''=-a$ [holds because $a$ has the form $C_1\cos\theta+C_2\sin\theta$],
\[
\begin{aligned}
r&=NR+\tfrac12aRR+\left(\tfrac13b-\frac{a'a'}{24N}\right)R^3+\left(\tfrac14c-\frac{a'b'}{24N}+\frac{aa'a'}{48NN}\right)R^4\ldots;\\
n^*\cos\varphi&=N+aR+\left(b-\frac{a'a'}{8N}\right)RR+\left(c-\frac{a'b'}{6N}+\frac{aa'a'}{12NN}\right)R^3\ldots\\
-n^*\sin\varphi&=\tfrac12a'R+\left(\tfrac13b'+\frac{aa'}{12N}\right)RR
+\left(\tfrac14c'+\frac{ab'}{24N}-\frac{a'b''}{24N}+\frac{{a'}^3}{48NN}-\frac{aaa'}{24NN}\right)R^3\ldots.
\end{aligned}
\]
Thus
\[
\begin{aligned}
-\sin\varphi&=\frac{a'}{2N}R+\left(\frac{b'}{3N}-\frac{5aa'}{12NN}\right)RR\ldots\\
-\varphi&=\frac{a'}{2N}R+\left(\frac{b'}{3N}-\frac{5aa'}{12NN}\right)RR\ldots\\
&=\tfrac12\alpha'R+\left(\tfrac13\beta'-\tfrac1{12}\alpha\alpha'\right)RR\ldots.
\end{aligned}
\]
The above partial differential equation further gives
\[
\begin{aligned}
-\psi&=\frac{a'}{2N}R+\left(\frac{b'}{6N}-\frac{aa'}{12NN}\right)RR\ldots\\
&=\tfrac12\alpha'R+\left(\tfrac16\beta'+\tfrac1{12}\alpha\alpha'\right)RR\ldots.
\end{aligned}
\]
If the positions of the points in the plane are expressed by $x$, $y$, where $x=R\cos\theta$, $y=R\sin\theta$, let
\par\endgroup
\clearpage
\noindent\mbox{}\hfill\textsc{Conformal Mapping of the Spheroid in the Plane.}\hfill165\par\medskip
\begingroup\small
\[
\log n^*=\log N-Ax-By-\tfrac12Cxx-Dxy-\tfrac12Eyy\ldots,
\]
from which we obtain
\[
\begin{array}{rl@{\qquad}rl}
\alpha R&=-Ax-By&\alpha'R&=Ay-Bx\\
\beta RR&=-\tfrac12Cxx-Dxy-\tfrac12Eyy&\beta'RR&=(C-E)xy-D(xx-yy)\\
&\text{etc.}&&\text{etc.}
\end{array}
\]
\[
\begin{aligned}
\varphi&=+\tfrac12Bx+(\tfrac13D+\tfrac1{12}AB)(xx-yy)\\
&\qquad-\tfrac12Ay+(-\tfrac13C+\tfrac13E-\tfrac1{12}AA+\tfrac1{12}BB)xy\ldots\\
\psi&=+\tfrac12Bx+(\tfrac16D-\tfrac1{12}AB)(xx-yy)\\
&\qquad-\tfrac12Ay+(-\tfrac16C+\tfrac16E+\tfrac1{12}AA-\tfrac1{12}BB)xy\ldots;\\[5pt]
r&=N\left(R+\tfrac12\alpha RR+(\tfrac13\beta+\tfrac16\alpha\alpha-\tfrac1{24}\alpha'\alpha')R^3\ldots\right)\\
\log r&=\log NR+\tfrac12\alpha R+\left(\tfrac13\beta+\tfrac1{24}(\alpha\alpha-\alpha'\alpha')\right)RR\ldots\\
&=\log N^0R+\left(\tfrac1{12}\beta+\tfrac1{24}(\alpha\alpha-\alpha'\alpha')\right)RR\ldots\\
&=\log N^0R-\tfrac1{24}(C+BB-AA)xx\\
&\qquad\qquad-\tfrac1{12}(D-2AB)xy\\
&\qquad\qquad-\tfrac1{24}(E+AA-BB)yy
\end{aligned}
\]
\begin{center}etc.,\end{center}
where $N^0$ is the value of $n^*$ at the midpoint of the straight line [so $\log N^0=\log N+\tfrac12\alpha R+\tfrac14\beta RR+\ldots$].

Now let
\[
\begin{aligned}
\mathopen{[}\log n=]\log\frac1{n^*}&=(e+e'x+e''xx+\ldots)(Y+y)^2\\
&\quad+(f+f'x+f''xx+\ldots)(Y+y)^4\\
&\quad+(g+g'x+g''xx+\ldots)(Y+y)^6
\end{aligned}
\]
\begin{center}etc.,\end{center}
with the second coordinate [at the starting point] [being] $=Y$. Then
\[
\begin{aligned}
-\log N&=eYY+fY^4+gY^6\ldots\\
A&=e'YY+f'Y^4+g'Y^6\ldots\\
B&=2eY+4fY^3+6gY^5\ldots\\
\tfrac12C&=e''YY+f''Y^4+g''Y^6\ldots\\
D&=2e'Y+4f'Y^3+6g'Y^5\ldots\\
\tfrac12E&=e+6fYY+15gY^4\ldots
\end{aligned}
\]
\begin{center}etc.\end{center}
\par\endgroup
\clearpage
\noindent166\hfill\textsc{Posthumous Papers.}\hfill\mbox{}\par\medskip
In the mapping of the spherical surface of radius $h$, $n^*$ is independent of $x$, and
\[
\begin{aligned}
\frac1{n^*}&=1+\tfrac12\frac{(Y+y)^2}{hh}+\tfrac1{24}\frac{(Y+y)^4}{h^4}\ldots\\
\log\frac1{n^*}&=\tfrac12\frac{(Y+y)^2}{hh}-\tfrac1{12}\frac{(Y+y)^4}{h^4}\ldots.
\end{aligned}
\]
Thus [we have]
\[
A=C=D=0
\]
[and]
\[
B=\frac{Y}{hh}-\tfrac13\frac{Y^3}{h^4}\ldots.
\]
The term $\tfrac1{12}BBxy$ [in $\varphi$ and $\psi$] therefore cannot exceed $\tfrac1{24}\frac{YY}{h^4}RR$, hence for $R=\lambda.100000\,\mathrm m$, $Y=\mu.400000\,\mathrm m$ [it cannot exceed] $\lambda\lambda.\mu\mu.0''{,}008$.

For $\varphi$ and $\psi$ the terms are [essentially, on the spheroid as well] accounted for by the formula $\tfrac12Bx-\tfrac12Ay$ if, in the two cases, the values of $B$ and $A$ are taken at points on the straight line lying $\frac23$ and $\frac13$ of its length from the starting point, respectively. These values are given above [Art.~15] to sufficient accuracy.

[If $A_{\frac23}$, $B_{\frac23}$ and $A_{\frac13}$, $B_{\frac13}$ are the values of $A$ and $B$ just described, then
\[
\begin{aligned}
\varphi&=\tfrac12B_{\frac23}.x-\tfrac12A_{\frac23}.y+\tfrac13(Dx-Cy)x+\tfrac1{12}(Ax+By)(Bx-Ay)\ldots\\
\psi&=\tfrac12B_{\frac13}.x-\tfrac12A_{\frac13}.y+\tfrac16(Dx-Cy)x+\tfrac1{12}(Ax+By)(-Bx+Ay)\ldots.]
\end{aligned}
\]
\begin{center}[17.]\end{center}
Generalizing the above principles leads to the following treatment.

Let $p$, $q$ be two variable quantities that determine the various points of the curved surface. Let $r$ denote the length of the shortest line and $\psi$ its direction at the start. Finally, let $\int\rho\,\mathrm d\psi$ be the length of the line of constant $r$ and $\mathrm ds$ [the] element of a line on the curved surface, expressed in general by $\sqrt{(A\,\mathrm dp^2+2B\,\mathrm dp\mathrm dq+C\,\mathrm dq^2)}$.

Set
\[
\begin{aligned}
\mathrm dr&=g\,\mathrm dp+h\,\mathrm dq\\
\rho\,\mathrm d\psi&=G\,\mathrm dp+H\,\mathrm dq.
\end{aligned}
\]
\clearpage
\noindent\mbox{}\hfill\textsc{Conformal Mapping of the Spheroid in the Plane.}\hfill167\par\medskip
From this [it follows that]
\[
\begin{aligned}
gg+GG&=A\\
gh+GH&=B\\
hh+HH&=C,
\end{aligned}
\]
thus [we obtain]
\[
(A-gg)(C-hh)=(B-gh)^2
\]
[or]
\[
\begin{array}{r@{\quad}l}
\mathrm I.&\begin{aligned}[t]
AC-BB&=Ahh-2Bgh+Cgg\\
&\left[=A\left(\frac{\partial r}{\partial q}\right)^2-2B\frac{\partial r}{\partial p}\frac{\partial r}{\partial q}+C\left(\frac{\partial r}{\partial p}\right)^2.\right.
\end{aligned}
\end{array}
\]
Furthermore, we have]
\[
\begin{aligned}
Ah-Bg&=G(Gh-Hg)\\
Bh-Cg&=H(Gh-Hg),
\end{aligned}
\]
thus
\[
\frac G\rho(Cg-Bh)+\frac H\rho(Ah-Bg)=0,
\]
i.e.
\[
\begin{array}{r@{\qquad}l}
\mathrm{II}.&\displaystyle\frac{\partial\psi}{\partial p}(Cg-Bh)+\frac{\partial\psi}{\partial q}(Ah-Bg)=0.
\end{array}
\]
If the angles that the lines of constant $p$ or $q$ make with the lines of constant $\psi$ are $\mathrm M$, $\mathrm N$, respectively, then
\[
\begin{array}{rl@{\qquad}rl}
g&=\cos\mathrm N\sqrt A&[G&=\sin\mathrm N\sqrt A\\
h&=\cos\mathrm M\sqrt C&H&=\sin\mathrm M\sqrt C;
\end{array}
\]
\rlap{\hspace{1.2em}\raisebox{-1.1ex}{.}}hence]
\[
B=\cos(\mathrm M-\mathrm N)\sqrt{AC}.
\]
If, for example,
\[
\begin{array}{rl@{\qquad}rl}
p&=R&x&=R\cos\theta\\
q&=\theta&y&=R\sin\theta\\
&&z&=1+tRR+uR^3\ldots
\end{array}
\]
where $t$, $u$ are functions of $\theta$, [then]
\[
\mathrm ds=\sqrt{\{\mathrm dR^2+RR\,\mathrm d\theta^2+\mathrm dz^2\}},
\]
and
\[
A=1+\left(\frac{\partial z}{\partial R}\right)^2,\qquad
B=\frac{\partial z}{\partial R}.\frac{\partial z}{\partial\theta},\qquad
C=RR+\left(\frac{\partial z}{\partial\theta}\right)^2.
\]
[By equation I, if $t'$, $u'$, $\ldots$ are the derivatives of $t$, $u$, $\ldots$ with respect to $\theta$, then
\clearpage
\noindent168\hfill\textsc{Posthumous Papers.}\hfill\mbox{}\par\medskip
\begingroup\small
\[
\begin{aligned}
&(1+t't'RR+2t'u'R^3\ldots)\left(\frac{\partial r}{\partial R}\right)^2
+(1+4ttRR+12tuR^3\ldots)\left(\frac{\partial r}{R\partial\theta}\right)^2\\
&\qquad-2(2tt'RR+(2tu'+3ut')R^3\ldots)\frac{\partial r}{\partial R}.\frac{\partial r}{R\partial\theta}\\
&\quad=1+(4tt+t't')RR+(12tu+2t'u')R^3\ldots,
\end{aligned}
\]
from which it follows that]
\[
r=R\qquad\ast\qquad+\tfrac23ttR^3+\tfrac32tuR^4\ldots.
\]
[Equation II then gives
\[
\begin{aligned}
&(1+(2tt+t't')RR+(6tu+2t'u')R^3\ldots)\frac{\partial\psi}{\partial R}\\
&+(-\tfrac23tt'RR-(\tfrac12tu'+\tfrac32ut')R^3\ldots)\frac{\partial\psi}{R\partial\theta}=0,
\end{aligned}
\]
consequently,
\[
\psi=\theta+\tfrac13tt'RR+(\tfrac12ut'+\tfrac16tu')R^3\ldots.]
\]
\bigskip
\begin{center}
On the Transformation of Coordinates.\par\medskip
[18.]
\end{center}
The fundamental equation is
\[
\mathrm f(x+iy+\omega)-\tfrac12\mathrm f(x+iy)+\tfrac12\mathrm f(x-iy)=\mathrm f(x+\Omega).
\]
We write
\[
\begin{aligned}
\mathrm f'(x)&=a\\
\mathrm f''(x)&=b\\
\mathrm f'''(x)&=c\\
\mathrm f^{\mathrm{IV}}(x)&=d
\end{aligned}
\]
\begin{center}etc.\end{center}
[Then we obtain]
\[
\begin{aligned}
&\mathrm f(x)+a\omega+\tfrac12b(\omega\omega+2iy\omega-yy)+\tfrac16c(\omega^3+3iy\omega\omega-3yy\omega)\\
&\qquad+\tfrac1{24}d(\omega^4+4iy\omega^3-6yy\omega\omega-4iy^3\omega+y^4)\ldots\\
&\quad=\mathrm f(x)+a\Omega+\tfrac12b\Omega\Omega+\tfrac16c\Omega^3+\tfrac1{24}d\Omega^4\ldots
\end{aligned}
\]
[and hence]
\par\endgroup
\clearpage
\noindent\mbox{}\hfill\textsc{Conformal Mapping of the Spheroid in the Plane.}\hfill169\par\medskip
\begingroup\small
\[
\begin{aligned}
x+\Omega&=x-\frac b{2a}yy+\left(\frac d{24a}\mathbin{\overset{\scriptscriptstyle .}{-}}\frac{b^3}{8a^3}\right)y^4\ldots\\
&\quad+\omega\left\{1+\frac ba.iy-\left(\frac c{2a}-\frac{bb}{2aa}\right)yy-\left(\frac d{6a}-\frac{b^3}{2a^3}\right)iy^3\ldots\right\}\\
&\quad+\omega\omega\left\{\left(\frac c{2a}-\frac{bb}{aa}\right)iy-\left(\frac d{4a}-\frac{3bc}{4aa}\right)yy\ldots\right\}\\
&\quad+\omega^3\left\{\left(\frac d{6a}-\frac{bc}{aa}+\frac{b^3}{a^3}\right)iy\ldots\right\}
\end{aligned}
\]
\begin{center}etc.\end{center}
\bigskip
\begin{center}[19.]\end{center}
The expansion is based on the following problem.

If
\[
u-Auu-Bu^3-Cu^4-Du^5-\ldots=p+Aq+Br+Cs+Dt+\ldots,
\]
[where $p$ is of the same order as $u$, $q$ is of order $uu$, etc.,] then
\[
\begin{aligned}
u&=p+A(pp+q)+AA(2p^3+2pq)+A^3(5p^4+6ppq+qq)\\
&\qquad\qquad+B(\ p^3+r)\qquad+AB(5p^4+3ppq+2pr)\\
&\qquad\qquad\qquad\qquad\qquad+C(\ p^4+s)\\
&\quad+\phantom{AA}A^4(14p^5+20p^3q+6pqq)\\
&\quad+AAB(21p^5+20p^3q+6ppr+3pqq+2qr)\\
&\quad+\phantom{A}AC(\ 6p^5+\ 4p^3q+2ps)\\
&\quad+\phantom{A}BB(\ 3p^5+3ppr)\\
&\quad+\phantom{AA}D(\ p^5+t)
\end{aligned}
\]
\begin{center}etc.\end{center}
In general,
\[
u=\sum A^\alpha B^\beta C^\gamma\ldots p^\pi q^\chi r^\rho s^\sigma\ldots.
\frac{\Pi(\alpha+\beta+\gamma+\ldots+\pi+\chi+\rho+\sigma+\ldots-1)}{\Pi(\alpha-\chi).\Pi(\beta-\rho).\Pi(\gamma-\sigma)\ldots\Pi\pi.\Pi\chi.\Pi\rho.\Pi\sigma\ldots},
\]
where
\[
\alpha,\ \beta,\ \gamma,\ \ldots,\ \pi,\ \chi,\ \rho,\ \sigma,\ \ldots,\ \alpha-\chi,\ \beta-\rho,\ \gamma-\sigma,\ \ldots
\]
must be nonnegative integers, and
\[
1+\alpha+2\beta+3\gamma+\ldots=\pi+2\chi+3\rho+4\sigma+\ldots
\]
[holds].
\par\endgroup
\vfill\noindent{\footnotesize IX.}\hfill22\hspace*{2cm}\par
\clearpage
\noindent170\hfill\textsc{Posthumous Papers.}\hfill\mbox{}\par\medskip
We now have
\[
u=\Omega
\]
[and]
\[
\begin{aligned}
\rlap{\hspace{-1.3em}.}p&=-z+(\omega+z)=\omega\\
q&=\phantom{z^3}-(\omega+z)^2\\
r&=z^3-(\omega+z)^3\\
s&=\phantom{z^3}-(\omega+z)^4\\
t&=z^5-(\omega+z)^5
\end{aligned}
\]
\begin{center}etc.\end{center}
\noindent\llap{\raisebox{-.4ex}{.}\,}[to set]. Thus
\begingroup\small
\[
\begin{aligned}
\Omega&=\omega-A(2\omega z+zz)-AA(4\omega\omega z+2\omega zz)\\
&\qquad\qquad\qquad\qquad-B(3\omega\omega z+3\omega zz)\\[5pt]
&\quad-A^3(\ 8\omega^3z\qquad\ast\qquad-4\omega z^3-z^4)\\
&\quad-AB(12\omega^3z+9\omega\omega zz\qquad\ast\qquad\ast\ )\\
&\quad-C(\ 4\omega^3z+6\omega\omega zz+4\omega z^3+z^4)\\[5pt]
&\qquad\qquad-A^4(16\omega^4z-16\omega^3zz-24\omega\omega z^3-6\omega z^4)\\
&\qquad\qquad-AAB(36\omega^4z\qquad\ast\qquad-30\omega\omega z^3-9\omega z^4)\\
&\qquad\qquad-AC(16\omega^4z+16\omega^3zz+8\omega\omega z^3+2\omega z^4)\\
&\qquad\qquad-BB(\ 9\omega^4z+9\omega^3zz\qquad\ast\qquad\ast\ )\\
&\qquad\qquad-D(\ 5\omega^4z+10\omega^3zz+10\omega\omega z^3+5\omega z^4)
\end{aligned}
\]
\begin{center}etc.\end{center}
\par\endgroup
To return to the previous notation [Art.~18], it only remains to set
\[
\begin{aligned}
z&=iy\\[5pt]
A&=-\frac b{2a}\\[5pt]
B&=-\frac c{6a}\\[5pt]
C&=-\frac d{24a}\\[5pt]
D&=-\frac e{120a}
\end{aligned}
\]
\begin{center}etc.\end{center}
\clearpage
\noindent\mbox{}\hfill\textsc{Conformal Mapping of the Spheroid in the Plane.}\hfill171\par\medskip
\begin{center}[20.]\end{center}
On the sphere [of radius $1$, according to p.~144, if $x$ is reckoned positively southward from an arbitrary starting point, we have
\[
\frac{\mathrm d\varphi}{\mathrm dx}=-1,\qquad\mathrm f'(x)=-\frac1{\cos\varphi};
\]
if one sets
\[
\sin\varphi=s,\qquad\cos\varphi=c,]\qquad\operatorname{tang}\varphi=t,
\]
[one therefore obtains
\begingroup\small
\[
\begin{array}{ll@{\quad}l}
\displaystyle\rlap{\hspace{-.3em}\raisebox{1.2ex}{.}}\mathrm f''(x)=+\frac{s}{cc}&\text{and hence}&A=+\tfrac12t\\[6pt]
\displaystyle\mathrm f'''(x)=-\frac{1+ss}{c^3}&&B=-\tfrac16(1+2tt)\\[6pt]
\displaystyle\mathrm f^{\mathrm{IV}}(x)=+\frac{5s+s^3}{c^4}&&C=+\tfrac1{24}(5t+6t^3)\\[6pt]
\displaystyle\mathrm f^{\mathrm V}(x)=-\frac{5+18ss+s^4}{c^5}&&D=-\tfrac1{120}(5+28tt+24t^4)\\[6pt]
\displaystyle\mathrm f^{\mathrm{VI}}(x)=+\frac{61s+58s^3+s^5}{c^6}&&E=+\tfrac1{720}(61t+180t^3+120t^5)\\[6pt]
\displaystyle\mathrm f^{\mathrm{VII}}(x)=-\frac{61+479ss+179s^4+s^6}{c^7}\Big]&&F=-\tfrac1{5040}(61+662tt+1320t^4+720t^6)
\end{array}
\]
\begin{center}etc.;\end{center}
\par\endgroup
[consequently, by Art.~19 or 20,]
\begingroup\footnotesize
\[
\begin{aligned}
x+\Omega&=x+\tfrac12tyy-(\tfrac5{24}t+\tfrac18t^3)y^4\ldots\\
&\quad+\omega\left\{1-\tfrac12(1+tt)yy+(\tfrac5{24}+\tfrac7{12}tt+\tfrac38t^4)y^4\ldots-i\left(ty-(\tfrac56t+\tfrac12t^3)y^3\ldots\right)\right\}\\
&\quad+\omega\omega\left\{\tfrac12tyy\ldots+i\left(\tfrac12y-(\tfrac5{12}+\tfrac14tt)y^3\ldots\right)\right\}\\
&\quad+\omega^3\left\{-(\tfrac16-\tfrac13tt)yy\ldots+i(\tfrac16ty\ldots)\right\}\\
&\quad+\omega^4\left\{\ldots-i(\tfrac1{24}y\ldots)\right\}
\end{aligned}
\]
\begin{center}etc.\end{center}
\par\endgroup
\bigskip
\begin{center}
[Series Relating $\varphi$, $\psi$ and $\omega$.]\par\medskip
[21.]
\end{center}
[For the numerical computation of $\varphi$ from $x$, Art.~1, or of $X$ from $\Phi$, Art.~13, the following equations between $\varphi$ and $\psi$ are used.

For the elliptical meridian arc with semimajor axis $a$ and eccentricity $e$, let the length from the equator to latitude $\varphi$ be $x$. Then we have]
\vfill\hfill$22^*$\hspace*{2cm}\par
\endgroup
